% Options for packages loaded elsewhere
\PassOptionsToPackage{unicode}{hyperref}
\PassOptionsToPackage{hyphens}{url}
\PassOptionsToPackage{dvipsnames,svgnames,x11names}{xcolor}
%
\documentclass[
]{article}
\usepackage{amsmath,amssymb}
\usepackage{iftex}
\ifPDFTeX
  \usepackage[T1]{fontenc}
  \usepackage[utf8]{inputenc}
  \usepackage{textcomp} % provide euro and other symbols
\else % if luatex or xetex
  \usepackage{unicode-math} % this also loads fontspec
  \defaultfontfeatures{Scale=MatchLowercase}
  \defaultfontfeatures[\rmfamily]{Ligatures=TeX,Scale=1}
\fi
\usepackage{lmodern}
\ifPDFTeX\else
  % xetex/luatex font selection
  \setmainfont[]{DejaVu Serif}
  \setmonofont[]{DejaVu Sans Mono}
\fi
% Use upquote if available, for straight quotes in verbatim environments
\IfFileExists{upquote.sty}{\usepackage{upquote}}{}
\IfFileExists{microtype.sty}{% use microtype if available
  \usepackage[]{microtype}
  \UseMicrotypeSet[protrusion]{basicmath} % disable protrusion for tt fonts
}{}
\makeatletter
\@ifundefined{KOMAClassName}{% if non-KOMA class
  \IfFileExists{parskip.sty}{%
    \usepackage{parskip}
  }{% else
    \setlength{\parindent}{0pt}
    \setlength{\parskip}{6pt plus 2pt minus 1pt}}
}{% if KOMA class
  \KOMAoptions{parskip=half}}
\makeatother
\usepackage{xcolor}
\usepackage[margin=1in]{geometry}
\setlength{\emergencystretch}{3em} % prevent overfull lines
\providecommand{\tightlist}{%
  \setlength{\itemsep}{0pt}\setlength{\parskip}{0pt}}
\setcounter{secnumdepth}{-\maxdimen} % remove section numbering
\usepackage{newunicodechar}
\newunicodechar{★}{\ensuremath{\star}}
\newunicodechar{✓}{\ensuremath{\checkmark}}
\newunicodechar{√}{\ensuremath{\surd}}
\newunicodechar{≈}{\ensuremath{\approx}}
\ifLuaTeX
  \usepackage{selnolig}  % disable illegal ligatures
\fi
\IfFileExists{bookmark.sty}{\usepackage{bookmark}}{\usepackage{hyperref}}
\IfFileExists{xurl.sty}{\usepackage{xurl}}{} % add URL line breaks if available
\urlstyle{same}
\hypersetup{
  colorlinks=true,
  linkcolor={Maroon},
  filecolor={Maroon},
  citecolor={Blue},
  urlcolor={Blue},
  pdfcreator={LaTeX via pandoc}}

\author{}
\date{}

\begin{document}

\hypertarget{theory-track}{%
\section{Theory track}\label{theory-track}}

Each chapter follows the same structure:

\begin{enumerate}
\def\labelenumi{\arabic{enumi}.}
\tightlist
\item
  \textbf{Motivation} --- why does this matter?
\item
  \textbf{Definitions} --- precise statements of objects we'll
  manipulate.
\item
  \textbf{Derivations / theorems} --- the math, with proofs where short
  enough to be illuminating.
\item
  \textbf{Worked example} --- a tiny numerical case you can verify by
  hand.
\item
  \textbf{Exercises} --- numbered, with difficulty markers:

  \begin{itemize}
  \tightlist
  \item
    ★ --- routine, should take \textless{} 5 min
  \item
    ★★ --- requires real thought, 10--30 min
  \item
    ★★★ --- open-ended or a small proof, 30+ min
  \end{itemize}
\item
  \textbf{Before the notebook} --- a prediction you must write down
  before running code.
\end{enumerate}

Do exercises in a paper notebook or a scratch markdown file. Solutions
are in \texttt{theory/solutions/} --- resist until you've made a genuine
attempt.

\hypertarget{math-prerequisites}{%
\subsection{Math prerequisites}\label{math-prerequisites}}

\begin{itemize}
\tightlist
\item
  Linear algebra: matrix multiplication, inner products,
  eigendecomposition, norms.
\item
  Probability: discrete distributions, expectation, variance,
  independence.
\item
  Calculus: partial derivatives, chain rule, gradient of scalar w.r.t.
  vector.
\end{itemize}

If any of those feel shaky, start with
\texttt{00\_linear\_algebra\_primer.md}.

\end{document}
